\documentclass[11pt]{article}
\usepackage[margin=1in]{geometry}
\usepackage{setspace}
\usepackage{enumitem}
\usepackage{amsmath}
\usepackage{amssymb}
\usepackage{fancyhdr}
\IfFileExists{lastpage.sty}{\usepackage{lastpage}}{}
\usepackage{hyperref}

\setlength{\parindent}{0pt}
\setlength{\parskip}{0.6em}
\setlist[itemize]{leftmargin=*}
\setlist[enumerate]{leftmargin=*}

\newcommand{\answerlines}[1]{\vspace{#1\baselineskip}}

\pagestyle{fancy}
\fancyhf{}
\IfFileExists{lastpage.sty}{
  \fancyfoot[C]{Financial Data Science II\ \ \ Updated March 17, 2026\ \ \ Page \thepage\ of \pageref{LastPage}}
}{
  \fancyfoot[C]{Financial Data Science II\ \ \ Updated March 17, 2026\ \ \ Page \thepage}
}

\begin{document}

\begin{center}
{\Large \textbf{Part 24: Signal Construction and Long-Short Portfolios}}
\end{center}

Parts 22 and 23 focused on constructing and screening candidate predictors. The
next task is to map a selected feature set into a trading rule.

A feature is not yet a signal. A signal is a rule for ranking opportunities or
assigning positions. The same feature can produce very different portfolios
depending on how it is standardized, combined, and rebalanced.

In this Part we use a simple cross-sectional example to illustrate this
transition from features to positions.

\textbf{Exercise.} Give an example of a quantity that is useful as a feature
but not yet well defined as a trading signal. What additional choices are
needed before it becomes a portfolio rule?

\answerlines{6}

\newpage

\section*{Basic Terms}

A \textbf{signal} is a rule that maps observed data into a trading score or
position. A \textbf{rebalance date} is a date at which the signal is recomputed
and the portfolio is updated. The \textbf{long leg} consists of the assets
receiving positive weight. The \textbf{short leg} consists of the assets
receiving negative weight. \textbf{Turnover} measures how much the portfolio
weights change from one rebalance date to the next.

We also distinguish between:

\begin{enumerate}
  \item \textbf{Cross-sectional signals:} rank many assets at the same time.
  \item \textbf{Time-series signals:} evaluate one asset relative to its own history.
\end{enumerate}

The example below is cross-sectional.

\textbf{Exercise.} Why might a feature that works well in a regression problem
still perform poorly when translated into a long-short portfolio?

\answerlines{6}

\newpage

\section*{Why the Mapping Matters}

Once a feature has been selected, several additional choices remain.

\begin{enumerate}
  \item Should the feature be standardized within each date?
  \item Should multiple features be combined additively or in some other way?
  \item Should the strategy hold the strongest names only, or all names with varying weights?
  \item How often should the portfolio rebalance?
  \item How should turnover and implementation cost be handled?
\end{enumerate}

These are portfolio-design choices rather than feature-engineering choices, but
they materially affect the final result.

\section*{A Simple Workflow}

A practical sequence is the following.

\begin{enumerate}
  \item Choose a universe and a rebalance schedule.
  \item Compute feature values using only information available at the rebalance date.
  \item Standardize or rank the features cross-sectionally.
  \item Combine the selected features into a composite signal.
  \item Map the signal into long and short weights.
  \item Evaluate return, volatility, Sharpe ratio, and turnover.
\end{enumerate}

\newpage

\section*{An Illustrative Monthly Long-Short Signal}

For illustration we use a small universe of large US equities. This is a
convenience sample, not a survivorship-free universe. The point is to
illustrate the mechanics of signal formation rather than to make strong claims
about a production strategy.

A simple signal can be constructed from three ingredients:

\begin{enumerate}
  \item medium-horizon momentum
  \item shorter-horizon momentum
  \item volatility and drawdown penalties
\end{enumerate}

These ingredients may be ranked cross-sectionally and combined into a composite
score. The portfolio rule then buys the highest-ranked names and shorts the
lowest-ranked names.

\section*{Interpreting the Example}

Several caveats are important.

\begin{enumerate}
  \item The universe is small and survivorship-biased.
  \item No transaction cost model has been applied.
  \item The weights are intentionally simple.
  \item The signal weights are chosen for illustration rather than by a formal estimation procedure.
\end{enumerate}

Even so, the example highlights the main distinction between a feature and a
portfolio rule. A feature is a candidate piece of information. A signal
specifies how that information is translated into a ranking and then into
positions.

\newpage

\section*{Signals in the LOB Setting}

The same logic applies at shorter horizons. A model of next-price direction or
a short-horizon imbalance measure is not yet a trading rule. One still needs to
specify thresholds, inventory controls, latency assumptions, and a mapping from
model output to actual orders.

\section*{Exercises}

\textbf{Exercise 1.} Modify the signal so that the long and short legs each
contain five assets rather than three. How does turnover change?

\answerlines{6}

\textbf{Exercise 2.} Replace the equal-weight portfolio rule with weights
proportional to the distance of each score from the cross-sectional median.
What practical issues arise?

\answerlines{6}

\textbf{Exercise 3.} Suppose that the signal is computed using month-end prices
but implemented at the close of the same day. Why might this introduce bias?

\answerlines{6}

\textbf{Exercise 4.} Give an example of a feature that may be useful for stock
selection but difficult to convert into a stable long-short signal.

\answerlines{6}

\section*{References}

\begin{enumerate}
  \item Grinold and Kahn, \textit{Active Portfolio Management}.
  \item Jegadeesh and Titman (1993), \textit{Returns to Buying Winners and Selling Losers}.
  \item Cochrane, \textit{Asset Pricing}.
  \item Lopez de Prado (2018), \textit{Advances in Financial Machine Learning}.
\end{enumerate}

\end{document}
